\section{The square-scale product grouping, with its missing directions (NS-95)}
\paragraph{Classification and attempted estimate.}
This is an exact arithmetic restriction, a scoped failure of an observation
comparison for arbitrary residuals, and a finite test on the actual target.
It does not prove the proposed eventual estimate
$E_{N^2}\le(1-c)E_N$ for a fixed $c>0$. We attempted to promote the
NS-86 divisor-product construction to a free signed block variational
problem, hoping that the product identity would supply this lower gain.
The calculation below retains the missing directions rather than identifying
the restricted problem with the full NS-87 observation energy.

Use the fixed $q=2$ Hilbert space and arithmetic atoms of NS-87. For
$M=N^2$, write $\phi_n=(I-\Pi_N)a_n$, $N<n\le M$, with complete
positive definite Gram $H$. Put $h_n=\langle r_N,\phi_n\rangle$ and
$\Delta=h^TH^{-1}h=E_N-E_M$. The positive definiteness follows from
the finite coefficient recovery of NS-78/81. The target is fixed throughout.

\subsection{Two distinct product spaces}
Let
\[
 {\cal P}_N=\{ij:1\le i,j\le N\}.
\]
Direct products give the span of $\phi_d$ for $d\in{\cal P}_N$, $d>N$.
Repeated factorizations do not give new directions. This is a proper
subspace of the full block: no prime $p>N$ is a product of two indices
at most $N$. For example $N^2-1>N$ is also missing for every $N\ge2$:
if both factors were at most $N$ and the product exceeded $N(N-1)$,
both would have to equal $N$.

NS-86 uses a different construction, completing each product by its
multiples. Free all of its signed product weights, putting
\[
 b_n=\sum_{\substack{d\mid n\\d\in{\cal P}_N}}\beta_d,
 \qquad N<n\le M.
\]
The weights $\beta_d$ can be arbitrary reals: each product has at least
one factor pair, and allowing arbitrary signed pair weights realizes any
such collection. Define the integer matrix
\[
 L_{n,d}=\boldsymbol1_{d\mid n},\quad
 N<n\le M,\ d\in{\cal P}_N.
\]
The resulting projected space has coefficient range $\operatorname{ran}L$.
Remove dependent columns by exact rational elimination, giving a basis $B$.
All old coefficients may be refitted without changing these new coefficients.

The NS-86 choice is $\beta_d=\sum_{ij=d}c_ic_j$, so its damped,
old-space-refitted direction belongs to this free space. Freeing the
weights generally loses NS-86's particular defect-convolution-square
identity; that identity is not asserted for the optimized free lift.
The direct-product and divisor-completed spaces are not presumed nested.

\subsection{A restriction on coefficients at large primes}
For every prime $p>N$ and integer $m\ge1$ with $mp\le N^2$,
\begin{equation}
 b_{mp}=\sum_{d\mid m}\beta_d.
 \label{eq:ns95-prime-fiber}
\end{equation}
Indeed no $d\in{\cal P}_N$ is divisible by $p$. Thus a product seed
dividing $mp$ must divide $m$. Conversely $m<N$, so every divisor of $m$
is already in ${\cal P}_N$ through the pair $(1,d)$. In particular
$b_p=\beta_1$ for all primes $N<p\le N^2$. This restriction survives
arbitrary signed weights, damping and old-space refitting. It is a
necessary restriction, not a complete characterization of the range of $L$.

If the block has two such primes $p,q$, let $\ell=e_p-e_q$ in new-index
coordinates. Then $B^T\ell=0$. The actual arithmetic vector
\[
 v=\sum_{n=N+1}^{M}(H^{-1}\ell)_n\phi_n
\]
is nonzero, orthogonal to every free-lift direction, and has squared
norm $\ell^TH^{-1}\ell>0$. Its full block projection is itself.
Consequently a comparison asserting a positive fraction of full block
gain from these observations for \emph{every} residual orthogonal to
$V_N$ is false, even in the original arithmetic family. This $v$ is a
constructed test residual, not $r_N=(I-\Pi_N)\tau$. No failure of an
actual-target lower estimate follows. The proposition is conditional
only on the specified pair of primes, and such pairs are explicitly
listed in each finite test; no asymptotic prime-count input is used.

There is also no positivity supplied by the direct product pairing
matrix $K_{ij}=\langle r_N,a_{ij}\rangle$. For $N\ge4$, its principal
minor on indices $2,N$ is
\[
 \det\begin{pmatrix}0&h_{2N}\\h_{2N}&h_{N^2}\end{pmatrix}
       =-h_{2N}^2.
\]
The zero is the old normal equation at index $4$. Whenever $h_{2N}\ne0$
this matrix is indefinite; the verifier checks nonzero values on the
three tested sizes. This pairing matrix is not the positive Gram $H$,
nor the pairing matrix of the divisor-completed lift. No eventual sign
or nonvanishing assertion is made.

\subsection{Exact gain split and a lower bound on the discarded gain}
The best free-lift step has coefficients and gain
\[
 x_L=B(B^THB)^{-1}B^Th,\qquad
 \Delta_L=h^TB(B^THB)^{-1}B^Th.
\]
Set $g=h-Hx_L$. Then $B^Tg=0$, and, with all cross terms retained,
\begin{equation}
 \Delta=\Delta_L+g^TH^{-1}g.
 \label{eq:ns95-exact-split}
\end{equation}
This follows either from orthogonal projection or by expanding the
right side and using $x_L^THx_L=h^Tx_L=\Delta_L$.
Both terms are nonnegative, but this alone does not compare either
to $E_N$ from below. In particular one must not drop the second term
while calling the first the complete joint observation energy.

For any nonzero coefficient constraint $\ell$ with $B^T\ell=0$, let
$x_*=H^{-1}h$. Cauchy--Schwarz in the $H$ metric gives
\begin{equation}
 \Delta-\Delta_L
 =\|x_*-x_L\|_H^2
 \ \ge\ \frac{(\ell^Tx_*)^2}{\ell^TH^{-1}\ell}.
 \label{eq:ns95-loss}
\end{equation}
The prime-pair constraint gives a strictly positive finite lower bound
in every tested block. This lower bound concerns the \emph{discarded}
gain and does not help prove a lower bound on the retained gain. It
quantifies the difference between the two optimization problems.

Every gain above uses the complete physical Gram, including its tail.
For a finite-observation implementation, NS-87's joint tail argument
also applies to either restricted space. Its $B_T\le e_T$ estimate
and $T=N^6$ localization remain valid with old-space rows retained.
This does not make the physical cutoff $N^2$ sufficient and supplies
no lower bound for the signed interior observations.

\subsection{Outcome and remaining input}
The finite tests compare the prescribed NS-86 direction, the direct
product span, the free divisor-completed span and the complete block
at $N=4,8,16$. They use exact rational rank reduction, complete Arb
Grams at two precisions, a doubled analytic cutoff and an independent
physical-cell integration of the free-lift approximant. These are tests
of the proposed grouping, not new best finite approximation errors.
The full $E_{N^2}$ values were already computed in NS-86.

To obtain the suggested cofinal square-scale contraction from this
free lift still requires an independent estimate
\[
 h^TB(B^THB)^{-1}B^Th\ge c E_N
\]
on the actual optimized residual, eventually on a specified squaring
sequence. This is stronger than a bound for the full block and is not
asserted necessary. The exact divisor-product algebra and old normal
equations have not supplied it. No fixed $c$ is inferred from the
finite table. An estimate can instead retain the complementary term
in~\eqref{eq:ns95-exact-split}. Neither alternative is proved here.

\paragraph{What had to change.}
The fixed nonlinear product weights were replaced by freely optimized
signed weights and all old coefficients were refitted. This repairs
much of the finite loss, at the cost of a larger arithmetic inverse and
loss of the defect-square prescription. Products lying below $N^2$
do not imply that product observations span that whole block. The
universal all-residual comparison is refuted under the explicit prime-pair
hypothesis; the actual-target cofinal estimate stays open. The original
Weil floor, G2 and RH remain open, and manuscript v1.66 is unchanged.
